\documentclass[11pt,a4paper]{article}
\usepackage[utf8]{inputenc}
\usepackage[T1]{fontenc}
\usepackage[catalan]{babel}
\usepackage{lmodern}
\linespread{1.22}
\usepackage[a4paper,top=2.4cm,bottom=2.2cm,left=1.5cm,right=1.5cm,headheight=1.1cm,headsep=0.5cm,footskip=1cm]{geometry}
\usepackage{amsmath,amssymb,cancel}
\usepackage{icomma}
\usepackage{graphicx}
\usepackage[table]{xcolor}
\usepackage{tikz}\usetikzlibrary{calc}
\usepackage{pgfplots}\pgfplotsset{compat=1.17}
\usepackage{fontawesome5}
\usepackage{ccicons}
\usepackage{fancyhdr}
\usepackage{enumitem}
\usepackage{booktabs}
\usepackage{tabularx}
\usepackage{needspace}
\usepackage[most]{tcolorbox}
\usepackage{titlesec}
\usepackage[hidelinks]{hyperref}
\hypersetup{pdftitle={Apunts de teoria. Unitat 3: Vectors i aplicacions. Rectes, plans i feixos de plans. 2n de Batxillerat MII},pdfsubject={Matemàtiques II},pdfkeywords={vectors a l'espai, dependència lineal, bases, producte escalar, producte vectorial, producte mixt, equacions de la recta, equacions del pla, posicions relatives, feixos de plans}}

% colors de la web (FentMates, bloc de Geometria: taronja; es manté el nom fmgreen dels apunts de les altres unitats)
\definecolor{fmnavy}{HTML}{1B2748}
\definecolor{fmgreen}{HTML}{DC7A29}
\definecolor{fmgreenl}{HTML}{F79645}
\definecolor{fmyellow}{HTML}{FFC94D}
\definecolor{fmgold}{HTML}{E9A919}
\definecolor{fmaccent}{HTML}{FF7A50}
\definecolor{fmteal}{HTML}{2BB3A3}
\definecolor{fmdef}{HTML}{1D4F73}\definecolor{fmdefbg}{HTML}{EEF4F9}
\definecolor{fmprop}{HTML}{3B4252}\definecolor{fmpropbg}{HTML}{F3F4F7}
\definecolor{fmex}{HTML}{A3203A}\definecolor{fmexbg}{HTML}{FBF0F2}
\definecolor{fmmuted}{HTML}{5D6685}\definecolor{fmball0}{HTML}{E2483D}\definecolor{fmball1}{HTML}{2F6FD6}
\definecolor{fmsurface}{HTML}{EEF1F8}

\newcommand{\lt}{<}\newcommand{\gt}{>}
% fórmula en línia que no es parteix mai; si no cap en una línia, es posa centrada
\newcommand{\fmm}[1]{\setbox0=\hbox{$#1$}\ifdim\wd0>\linewidth\fmdisp{#1}\else\mbox{\box0}\fi}
% fórmula centrada; si és més ampla que la línia, s'encongeix fins que hi cap
\newcommand{\fmdisp}[1]{\[\setbox2=\hbox{$\displaystyle #1$}\ifdim\wd2>\linewidth\resizebox{\linewidth}{!}{\box2}\else\box2\fi\]}
\DeclareUnicodeCharacter{2212}{-}
\DeclareUnicodeCharacter{2225}{\ensuremath{\parallel}}\DeclareUnicodeCharacter{22A5}{\ensuremath{\perp}}\DeclareUnicodeCharacter{2226}{\ensuremath{\nparallel}}
\DeclareUnicodeCharacter{207A}{\ensuremath{^{+}}}\DeclareUnicodeCharacter{207B}{\ensuremath{^{-}}}
\setlength{\parindent}{0pt}
\setlength{\parskip}{12pt}
\setlength{\lineskip}{5pt}\setlength{\lineskiplimit}{5pt}  % més aire al voltant de les fraccions en línia
\setlength{\emergencystretch}{3em}
\setlist{itemsep=4pt,topsep=4pt}
\renewcommand{\thesection}{3.\arabic{section}}
\titleformat{\section}{\color{fmnavy}\Large\bfseries}{\colorbox{fmdef}{\color{white}\strut\ \thesection\ }}{0.6em}{}[{\color{fmdef}\titlerule[1.2pt]}]
\titlespacing*{\section}{0pt}{16pt}{8pt}
\titleformat*{\subsection}{\color{fmgreen}\large\bfseries}
\allowdisplaybreaks

% dibuix de la targeta de la unitat a la web (illustrations.js), en blanc
\newcommand{\fmcard}{\begin{tikzpicture}[x=0.062cm,y=-0.062cm,white,line width=3.2pt,line cap=round,line join=round]
  \begin{scope}[transparency group,opacity=0.15]
  \draw[line width=2.6pt,opacity=0.7] (110,70)--(110,8) (110,70)--(210,92) (110,70)--(40,108);
  \fill[opacity=0.14] (50,40)--(150,22)--(210,62)--(110,82)--cycle;
  \draw (50,40)--(150,22)--(210,62)--(110,82)--cycle;
  \draw (60,110)--(190,14);
  \fill (128,47) circle (5.5);
  \draw[line width=3.4pt] (128,47)--(128,18) (122,26)--(128,16)--(134,26);
  \node[anchor=base west,inner sep=0pt,font=\fontsize{22}{22}\selectfont] at (136,22) {$\vec n$};
  \node[anchor=base west,inner sep=0pt,font=\fontsize{22}{22}\selectfont] at (196,112) {$\pi$};
  \end{scope}
\end{tikzpicture}}
\newcommand{\fmbrand}{\raisebox{-0.25\height}{\includegraphics[height=0.5cm]{logo-fentmates.pdf}}}

% capçalera i peu
\pagestyle{fancy}\fancyhf{}
\fancyhead[L]{\small\fmbrand}
\fancyhead[R]{\small\color{fmmuted}2n Batxillerat MII · Unitat 3: Vectors, rectes i plans}
\fancyfoot[L]{\scriptsize\color{fmmuted}\href{https://creativecommons.org/licenses/by-nc-sa/4.0/deed.ca}{\ccbyncsa}\ \ FentMates · fentmates.org · CC BY-NC-SA 4.0}
\fancyfoot[R]{\small\color{fmmuted}\thepage}
\renewcommand{\headrulewidth}{0.4pt}\renewcommand{\headrule}{\color{fmmuted!40}\hrule height \headrulewidth}
\fancypagestyle{primera}{\fancyhf{}\renewcommand{\headrulewidth}{0pt}
  \fancyfoot[L]{\scriptsize\color{fmmuted}\href{https://creativecommons.org/licenses/by-nc-sa/4.0/deed.ca}{\ccbyncsa}\ \ FentMates · fentmates.org · CC BY-NC-SA 4.0}
  \fancyfoot[R]{\small\color{fmmuted}\thepage}}

% caixes (mateixos colors que la web)
% quan una caixa es parteix entre dues pàgines, la vora del tall és discontínua
\tcbset{fmbase/.style={enhanced,breakable,arc=3mm,boxrule=1pt,left=3.5mm,right=3.5mm,top=2mm,bottom=2.5mm,before skip=10pt,after skip=10pt,fonttitle=\bfseries,
  frame style={fill=tcbcolframe},
  overlay first={\fill[tcbcolback] ([xshift=1pt]frame.south west) rectangle ([xshift=-1pt,yshift=1.2pt]frame.south east);
    \draw[tcbcolframe,line width=1pt,dash pattern=on 5pt off 3pt] (frame.south west)--(frame.south east);},
  overlay middle={\fill[tcbcolback] ([xshift=1pt]frame.south west) rectangle ([xshift=-1pt,yshift=1.2pt]frame.south east);
    \fill[tcbcolback] ([xshift=1pt]frame.north west) rectangle ([xshift=-1pt,yshift=-1.2pt]frame.north east);
    \draw[tcbcolframe,line width=1pt,dash pattern=on 5pt off 3pt] (frame.south west)--(frame.south east) (frame.north west)--(frame.north east);},
  overlay last={\fill[tcbcolback] ([xshift=1pt]frame.north west) rectangle ([xshift=-1pt,yshift=-1.2pt]frame.north east);
    \draw[tcbcolframe,line width=1pt,dash pattern=on 5pt off 3pt] (frame.north west)--(frame.north east);}}}
\newtcolorbox{fmdef}[1]{fmbase,colback=fmdefbg,colframe=fmdef,coltitle=white,title={#1}}
\newtcolorbox{fmprop}[1]{fmbase,colback=fmpropbg,colframe=fmprop,coltitle=white,title={#1},halign title=center}
\newtcolorbox{fmex}[1]{fmbase,before upper=\raggedright,colback=fmexbg,colframe=fmex,colbacktitle=fmexbg,coltitle=fmex,title={\faPencil*\ #1},
  titlerule=0.4pt,titlerule style={fmex!40,dashed}}
\newtcolorbox{fmwarn}[1]{fmbase,colback=fmsf!6,colframe=fmsf,coltitle=white,title={#1}}
\newtcolorbox{fmnote}{enhanced,breakable,colback=white,colframe=white,borderline west={3pt}{0pt}{fmaccent},arc=0pt,left=3mm,right=2mm,top=1mm,bottom=1mm,before skip=6pt,after skip=6pt,before upper=\raggedright}
\newcommand{\fmkey}[1]{\colorbox{fmyellow!45}{\strut #1}}
\newcommand{\fmuline}[1]{\textbf{\color{fmaccent}#1}}
\definecolor{fmok}{HTML}{15803D}\definecolor{fmev}{HTML}{2563EB}\definecolor{fmsf}{HTML}{C2410C}\definecolor{fmsi}{HTML}{7C3AED}\definecolor{fmas}{HTML}{A3203A}
\newcommand{\fmbadge}[2]{\tcbox[on line,colback=#1,colframe=#1,arc=2mm,boxsep=0pt,left=4pt,right=4pt,top=1.5pt,bottom=1.5pt]{\color{white}\footnotesize\bfseries #2}}
\newcommand{\fmchip}[1]{\tcbox[on line,colback=fmsf!10,colframe=fmsf!50,arc=2mm,boxsep=1pt,left=4pt,right=4pt]{#1}}
\newcommand{\fmlevel}[2]{\tcbox[on line,colback=white,colframe=fmmuted!40,arc=2mm,boxsep=2pt,halign=center]{\begin{tabular}{c}{\scriptsize\color{fmmuted}\MakeUppercase{#1}}\\#2\end{tabular}}}
% gràfiques
\pgfplotsset{fmaxis/.style={axis lines=middle,grid=major,grid style={fmmuted!15},axis line style={fmmuted,-latex},
  tick label style={font=\scriptsize,color=fmmuted},clip=true,unbounded coords=jump,restrict y to domain=-200:200,
  every axis plot/.append style={line join=round}},
  fmcurve/.style={fmaccent,line width=1.6pt,smooth}}
\tikzset{fmasy/.style={fmmuted,dashed,line width=0.9pt},fmguide/.style={fmaccent,dotted,line width=0.8pt},
  fmdot/.style={circle,fill=fmaccent,inner sep=0pt,minimum size=5pt},fmopen/.style={circle,draw=fmaccent,fill=white,line width=1pt,inner sep=0pt,minimum size=5pt}}
\newcommand{\fmpau}[1]{\tcbox[on line,colback=fmnavy,colframe=fmnavy,arc=2mm,boxsep=0pt,left=4pt,right=4pt,top=1.5pt,bottom=1.5pt]{\color{white}\scriptsize\bfseries\MakeUppercase{#1}}\ }

\begin{document}
\thispagestyle{primera}
\begin{tikzpicture}[remember picture,overlay]
  \fill[fmgreen] (current page.north west) rectangle ([yshift=-10.1cm]current page.north east);
  \begin{scope}
    \clip (current page.north west) rectangle ([yshift=-10.1cm]current page.north east);
    \node[anchor=north west,inner sep=0pt] at ([xshift=6.6cm,yshift=-1.3cm]current page.north west) {\fmcard};
  \end{scope}
  \node[anchor=north west,inner sep=0pt] at ([xshift=1.5cm,yshift=-1.1cm]current page.north west) {\tcbox[colback=white,colframe=white,arc=2.5mm,boxsep=0pt,left=7pt,right=7pt,top=5pt,bottom=5pt]{\includegraphics[height=0.75cm]{logo-fentmates.pdf}}};
\end{tikzpicture}
\vspace*{0.0cm}
\tcbox[colback=white!15!fmgreen,colframe=white!15!fmgreen,arc=2.5mm,boxsep=0pt,left=6pt,right=6pt,top=3pt,bottom=3pt]{\color{white}\small\bfseries Apunts de teoria · 2n Batxillerat · Matemàtiques II}

\vspace{4pt}
{\raggedright\color{white}\fontsize{25}{30}\selectfont\bfseries{\color{fmyellow}Unitat 3} · Vectors i aplicacions. Rectes, plans i feixos de plans\par}
\vspace{6pt}
\begin{minipage}{\linewidth}\color{white!88!fmgreen}\small Teoria i exemples resolts pas a pas: vectors a l'espai i les seues aplicacions (punt mig, simètric, punts alineats, baricentre), dependència lineal i bases, els productes escalar, vectorial i mixt (angles, àrees i volums), totes les equacions de la recta i del pla, les posicions relatives de rectes i plans, la discussió de tres plans amb el teorema de Rouché i els feixos de plans.\end{minipage}\par
% l'índex comença sempre a la mateixa altura, per davall de la franja (com en les unitats 5-10)
\vspace*{\dimexpr 8.6cm-\pagetotal\relax}

\begin{tcolorbox}[colback=fmsurface,colframe=fmsurface,arc=3mm,left=4mm,right=4mm,top=2mm,bottom=2mm]
\small\textbf{Índex}\par\smallskip
\begin{minipage}[t]{0.49\linewidth}\small \mbox{\hyperlink{sec1}{\textbf{\color{fmgreen}3.1}\ \ Vectors a l'espai}}\\[3pt]
\mbox{\hyperlink{sec2}{\textbf{\color{fmgreen}3.2}\ \ Aplicacions dels vectors}}\\[3pt]
\mbox{\hyperlink{sec3}{\textbf{\color{fmgreen}3.3}\ \ Dependència lineal i bases}}\\[3pt]
\mbox{\hyperlink{sec4}{\textbf{\color{fmgreen}3.4}\ \ Producte escalar}}\\[3pt]
\mbox{\hyperlink{sec5}{\textbf{\color{fmgreen}3.5}\ \ Producte vectorial}}\\[3pt]
\mbox{\hyperlink{sec6}{\textbf{\color{fmgreen}3.6}\ \ Producte mixt}}\\[3pt]\end{minipage}\hfill\begin{minipage}[t]{0.49\linewidth}\small \mbox{\hyperlink{sec7}{\textbf{\color{fmgreen}3.7}\ \ Equacions de la recta}}\\[3pt]
\mbox{\hyperlink{sec8}{\textbf{\color{fmgreen}3.8}\ \ Equacions del pla}}\\[3pt]
\mbox{\hyperlink{sec9}{\textbf{\color{fmgreen}3.9}\ \ Posicions relatives de rectes i plans}}\\[3pt]
\mbox{\hyperlink{sec10}{\textbf{\color{fmgreen}3.10}\ \ Posició relativa de tres plans}}\\[3pt]
\mbox{\hyperlink{sec11}{\textbf{\color{fmgreen}3.11}\ \ Feixos de plans}}\\[3pt]\end{minipage}
\end{tcolorbox}

\Needspace{18\baselineskip}
\section{Vectors a l'espai}\hypertarget{sec1}{}

\Needspace{7\baselineskip}
\begin{fmdef}{Definicions}
\begin{itemize}
\item Un \textbf{sistema de referència} de l'espai està format per un punt \fmm{O} (l'\textbf{origen}) i tres vectors \fmm{\vec i=(1,0,0)}, \fmm{\vec j=(0,1,0)}, \fmm{\vec k=(0,0,1)}, que determinen els eixos \fmm{OX}, \fmm{OY} i \fmm{OZ}. Cada punt \fmm{P} té tres coordenades \fmm{P(x,y,z)}, que són les components del seu \textbf{vector de posició} \fmm{\overrightarrow{OP}=x\,\vec i+y\,\vec j+z\,\vec k}.
\item Les \textbf{components} del vector \fmm{\overrightarrow{AB}} són les coordenades de l'extrem menys les de l'origen:

\fmdisp{A(a_1,a_2,a_3),\ B(b_1,b_2,b_3)\ \Rightarrow\ \overrightarrow{AB}=(b_1-a_1,\ b_2-a_2,\ b_3-a_3)}

\item El \textbf{mòdul} és la longitud del vector:

\fmdisp{|\vec v|=\sqrt{v_1^2+v_2^2+v_3^2}\qquad\qquad d(A,B)=|\overrightarrow{AB}|=\sqrt{(b_1-a_1)^2+(b_2-a_2)^2+(b_3-a_3)^2}}

\item Com en el pla, dos vectors són \textbf{iguals} (equipo\lgem{}ents) si tenen el mateix mòdul, direcció i sentit, és a dir, les mateixes components. Un vector és \textbf{unitari} si té mòdul 1; per a \textbf{normalitzar} un vector el dividim pel seu mòdul: \fmm{\dfrac{\vec v}{|\vec v|}}.
\end{itemize}
\end{fmdef}

\Needspace{7\baselineskip}
\begin{fmprop}{Operacions amb vectors}
Si \fmm{\vec u=(u_1,u_2,u_3)}, \fmm{\vec v=(v_1,v_2,v_3)} i \fmm{k\in\mathbb R}, les operacions es fan component a component:

\fmdisp{\vec u+\vec v=(u_1+v_1,\ u_2+v_2,\ u_3+v_3)\qquad\qquad k\,\vec u=(k\,u_1,\ k\,u_2,\ k\,u_3)}

Gràficament, \fmm{\vec u+\vec v} és la diagonal del para\lgem{}elogram que formen \fmm{\vec u} i \fmm{\vec v}; \fmm{k\,\vec u} té la mateixa direcció que \fmm{\vec u}, el mateix sentit si \fmm{k\gt0} i sentit contrari si \fmm{k\lt0}. Les propietats (commutativa, associativa, distributives\dots{}) són les mateixes que en el pla.
\end{fmprop}

\Needspace{7\baselineskip}
\begin{fmex}{Exemple resolt}
\begin{enumerate}[label=\textbf{\arabic*.},leftmargin=*]
\item Siguen \fmm{A(2,-1,3)} i \fmm{B(5,1,-3)}. Calcula \fmm{\overrightarrow{AB}}, la distància entre \fmm{A} i \fmm{B} i un vector unitari amb la mateixa direcció i sentit que \fmm{\overrightarrow{AB}}.

\fmm{\overrightarrow{AB}=\big(5-2,\ 1-(-1),\ -3-3\big)=(3,\ 2,\ -6)}

\fmm{d(A,B)=|\overrightarrow{AB}|=\sqrt{3^2+2^2+(-6)^2}=\sqrt{9+4+36}=\sqrt{49}=7}

\fmm{\vec w=\dfrac{\overrightarrow{AB}}{|\overrightarrow{AB}|}=\left(\dfrac37,\ \dfrac27,\ -\dfrac67\right)}. Comprovació: \fmm{\dfrac{9+4+36}{49}=1\ \checkmark}
\item Si \fmm{\vec u=(1,-2,3)} i \fmm{\vec v=(2,0,-1)}, calcula \fmm{2\vec u-3\vec v}.

\fmm{2\vec u-3\vec v=(2,-4,6)-(6,0,-3)=(-4,\ -4,\ 9)}
\end{enumerate}
\end{fmex}

\Needspace{5\baselineskip}
\section{Aplicacions dels vectors}\hypertarget{sec2}{}

\Needspace{7\baselineskip}
\begin{fmdef}{Fórmules}
\begin{itemize}
\item \textbf{Punt mig} del segment \fmm{AB}:
          \fmm{\ M=\left(\dfrac{a_1+b_1}{2},\ \dfrac{a_2+b_2}{2},\ \dfrac{a_3+b_3}{2}\right)}
\item \textbf{Simètric} de \fmm{A} respecte d'un punt \fmm{M}: com que \fmm{M} és el punt mig de \fmm{A} i \fmm{A'},
          \fmm{\ A'=2M-A=(2m_1-a_1,\ 2m_2-a_2,\ 2m_3-a_3)}
\item \textbf{Dividir un segment en \fmm{n} parts iguals}: els punts de divisió són
          \fmm{\ P_k=A+\dfrac{k}{n}\,\overrightarrow{AB}}, \fmm{\ k=1,\dots,n-1}.
\item Tres punts \fmm{A}, \fmm{B}, \fmm{C} estan \textbf{alineats} si \fmm{\overrightarrow{AB}} i \fmm{\overrightarrow{AC}} són \textbf{proporcionals} (tenen la mateixa direcció):\\[2ex]
          \fmm{\ \dfrac{c_1-a_1}{b_1-a_1}=\dfrac{c_2-a_2}{b_2-a_2}=\dfrac{c_3-a_3}{b_3-a_3}} (equivalentment, \fmm{\mathrm{rg}\big(\overrightarrow{AB},\overrightarrow{AC}\big)=1}).
          \newpage
\item \textbf{Baricentre} d'un triangle \fmm{ABC} (punt de tall de les mitjanes):
          \fmm{\ G=\left(\dfrac{a_1+b_1+c_1}{3},\ \dfrac{a_2+b_2+c_2}{3},\ \dfrac{a_3+b_3+c_3}{3}\right)}
\item \textbf{Para\lgem{}elogram} \fmm{ABCD}: \fmm{\overrightarrow{AB}=\overrightarrow{DC}}, i per tant el quart vèrtex és \fmm{D=A+C-B}.
\end{itemize}
\end{fmdef}

\Needspace{7\baselineskip}
\begin{fmex}{Exemple resolt}
\begin{enumerate}[label=\textbf{\arabic*.},leftmargin=*]
\item Calcula el punt mig de \fmm{A(2,-1,3)} i \fmm{B(4,5,-1)}, i el simètric de \fmm{P(1,2,3)} respecte de \fmm{Q(2,0,-1)}.

\fmm{M=\left(\dfrac{2+4}{2},\ \dfrac{-1+5}{2},\ \dfrac{3+(-1)}{2}\right)=(3,\ 2,\ 1)}

\fmm{P'=2Q-P=(4-1,\ 0-2,\ -2-3)=(3,\ -2,\ -5)}. Comprovació: el punt mig de \fmm{P} i \fmm{P'} és \fmm{(2,0,-1)=Q\ \checkmark}
\item Estan alineats els punts \fmm{A(1,2,-1)}, \fmm{B(3,1,1)} i \fmm{C(7,-1,5)}?

\fmm{\overrightarrow{AB}=(2,-1,2)}, \fmm{\quad\overrightarrow{AC}=(6,-3,6)}

\fmm{\dfrac62=\dfrac{-3}{-1}=\dfrac66=3\Rightarrow\overrightarrow{AC}=3\,\overrightarrow{AB}}. Són proporcionals: els punts \textbf{estan alineats}.
\item Divideix el segment d'extrems \fmm{A(1,-2,4)} i \fmm{B(7,4,-2)} en tres parts iguals.

\fmm{\overrightarrow{AB}=(6,6,-6)\Rightarrow\dfrac13\overrightarrow{AB}=(2,2,-2)}

\fmm{P_1=A+\tfrac13\overrightarrow{AB}=(3,\ 0,\ 2)}, \fmm{\quad P_2=A+\tfrac23\overrightarrow{AB}=(5,\ 2,\ 0)}
\end{enumerate}
\end{fmex}

\Needspace{18\baselineskip}
\section{Dependència lineal i bases}\hypertarget{sec3}{}

\Needspace{7\baselineskip}
\begin{fmdef}{Definicions}
\begin{itemize}
\item Un vector \fmm{\vec w} és \textbf{combinació lineal} de \fmm{\vec u} i \fmm{\vec v} si existeixen nombres \fmm{a,b} tals que \fmm{\vec w=a\,\vec u+b\,\vec v}.
\item Uns vectors són \textbf{linealment dependents} si algun és combinació lineal dels altres, i \textbf{linealment independents} en cas contrari. Geomètricament: dos vectors són dependents si són \textbf{para\lgem{}els} (proporcionals); tres vectors són dependents si són \textbf{coplanaris} (es poden dibuixar en un mateix pla).
\item El \textbf{rang} d'un conjunt de vectors és el nombre màxim de vectors linealment independents que hi ha. Es calcula com el rang de la matriu que té els vectors per files (unitat 1).
\item Una \textbf{base} de l'espai és qualsevol conjunt de \textbf{tres vectors linealment independents}. Llavors tot vector s'escriu de manera única com a combinació lineal d'ells: \fmm{\vec x=a\,\vec u+b\,\vec v+c\,\vec w}; els nombres \fmm{(a,b,c)} són les \textbf{coordenades} de \fmm{\vec x} en aquesta base. La base \fmm{\{\vec i,\vec j,\vec k\}} és la \textbf{base canònica}.
\end{itemize}
\end{fmdef}

\Needspace{7\baselineskip}
\begin{fmprop}{Criteri del determinant}

\fmdisp{\{\vec u,\vec v,\vec w\}\ \text{és base} \iff\begin{vmatrix}u_1&u_2&u_3\\v_1&v_2&v_3\\w_1&w_2&w_3\end{vmatrix}\neq0}

Si el determinant és 0, els tres vectors són coplanaris (rang 2) o, fins i tot, para\lgem{}els (rang 1). En l'espai, \textbf{quatre o més vectors sempre són linealment dependents}.
\end{fmprop}

\Needspace{7\baselineskip}
\begin{fmex}{Exemple resolt}
\begin{enumerate}[label=\textbf{\arabic*.},leftmargin=*]
\item Estudia si \fmm{\vec u=(1,2,-1)}, \fmm{\vec v=(0,1,3)} i \fmm{\vec w=(2,5,1)} formen una base. Si no en formen, escriu \fmm{\vec w} com a combinació lineal de \fmm{\vec u} i \fmm{\vec v}.

\fmm{\begin{vmatrix}1&2&-1\\0&1&3\\2&5&1\end{vmatrix}=1\cdot(1-15)-2\cdot(0-6)+(-1)\cdot(0-2)=-14+12+2=0}

El determinant és 0: \textbf{no formen una base} (són coplanaris). Com que \fmm{\vec u} i \fmm{\vec v} no són proporcionals, el rang és 2.

\fmm{(2,5,1)=a(1,2,-1)+b(0,1,3)\Rightarrow\begin{cases}a=2\\2a+b=5\\-a+3b=1\end{cases}\Rightarrow a=2,\ b=1}. La tercera equació també es compleix: \fmm{-2+3=1\ \checkmark}

\fmm{\vec w=2\,\vec u+\vec v}
\item Per a quins valors de \fmm{m} formen una base els vectors \fmm{(1,1,m)}, \fmm{(m,1,1)} i \fmm{(1,m,1)}?

\fmm{\begin{vmatrix}1&1&m\\m&1&1\\1&m&1\end{vmatrix}=(1-m)-(m-1)+m(m^2-1)=(m-1)(m^2+m-2)=(m-1)^2(m+2)}

S'anu\lgem{}a per a \fmm{m=1} i \fmm{m=-2}. Per tant, \textbf{formen base si \fmm{m\neq1} i \fmm{m\neq-2}}.
\end{enumerate}
\end{fmex}

\Needspace{7\baselineskip}
\section{Producte escalar}\hypertarget{sec4}{}

\Needspace{7\baselineskip}
\begin{fmdef}{Definició i expressió analítica}
El \textbf{producte escalar} de dos vectors és un \textbf{nombre}:

\fmdisp{\vec u\cdot\vec v=|\vec u|\,|\vec v|\cos\alpha\qquad\qquad \vec u\cdot\vec v=u_1v_1+u_2v_2+u_3v_3}

on \fmm{\alpha} és l'angle que formen (en una base ortonormal com la canònica, les dues expressions coincideixen).

\begin{itemize}
\item \textbf{Propietats}: \fmm{\vec u\cdot\vec v=\vec v\cdot\vec u}; \fmm{\ \vec u\cdot(\vec v+\vec w)=\vec u\cdot\vec v+\vec u\cdot\vec w}; \fmm{\ (k\vec u)\cdot\vec v=k(\vec u\cdot\vec v)}; \fmm{\ \vec u\cdot\vec u=|\vec u|^2}.
\item \textbf{Angle entre dos vectors}: \fmm{\ \cos\alpha=\dfrac{\vec u\cdot\vec v}{|\vec u|\,|\vec v|}}, \fmm{\ 0^\circ\le\alpha\le180^\circ}.
\item \textbf{Perpendicularitat} (ortogonalitat): \fmm{\ \vec u\perp\vec v\iff\vec u\cdot\vec v=0}.
\item \textbf{Projecció} de \fmm{\vec u} sobre \fmm{\vec v}: la projecció escalar és \fmm{\dfrac{\vec u\cdot\vec v}{|\vec v|}} (amb signe) i el vector projecció és \fmm{\dfrac{\vec u\cdot\vec v}{|\vec v|^2}\,\vec v}.
\item \textbf{Cosinus directors}: els cosinus dels angles que forma \fmm{\vec v} amb els eixos són \fmm{\cos\alpha=\dfrac{v_1}{|\vec v|}}, \fmm{\cos\beta=\dfrac{v_2}{|\vec v|}}, \fmm{\cos\gamma=\dfrac{v_3}{|\vec v|}}, i compleixen \fmm{\cos^2\alpha+\cos^2\beta+\cos^2\gamma=1}.
\end{itemize}
\end{fmdef}

\Needspace{7\baselineskip}
\begin{fmex}{Exemple resolt}
\begin{enumerate}[label=\textbf{\arabic*.},leftmargin=*]
\item Calcula el producte escalar i l'angle de \fmm{\vec u=(2,-1,3)} i \fmm{\vec v=(1,4,-2)}.

\fmm{\vec u\cdot\vec v=2\cdot1+(-1)\cdot4+3\cdot(-2)=2-4-6=-8}

\fmm{|\vec u|=\sqrt{4+1+9}=\sqrt{14}}, \fmm{\quad|\vec v|=\sqrt{1+16+4}=\sqrt{21}}

\fmm{\cos\alpha=\dfrac{-8}{\sqrt{14}\sqrt{21}}=\dfrac{-8}{7\sqrt6}\approx-0{,}4666\Rightarrow\alpha\approx117^\circ48'43''}. Com que el producte escalar és negatiu, l'angle és obtús.
\item Troba \fmm{k} perquè \fmm{\vec u=(k,2,-1)} i \fmm{\vec v=(3,k,4)} siguen perpendiculars.

\fmm{\vec u\cdot\vec v=3k+2k-4=5k-4=0\Rightarrow k=\dfrac45}
\end{enumerate}
\end{fmex}

\Needspace{18\baselineskip}
\section{Producte vectorial}\hypertarget{sec5}{}

\Needspace{7\baselineskip}
\begin{fmdef}{Definició}
El \textbf{producte vectorial} \fmm{\vec u\times\vec v} és un \textbf{vector}:

\begin{itemize}
\item de \textbf{mòdul} \fmm{|\vec u\times\vec v|=|\vec u|\,|\vec v|\sin\alpha};
\item de \textbf{direcció} perpendicular a \fmm{\vec u} i a \fmm{\vec v};
\item de \textbf{sentit} el d'avanç d'un tirabuixó que gira de \fmm{\vec u} cap a \fmm{\vec v} (regla de la mà dreta).
\end{itemize}

Es calcula amb un determinant «simbòlic» desenvolupat per la primera fila:

\fmdisp{\vec u\times\vec v=\begin{vmatrix}\vec i&\vec j&\vec k\\u_1&u_2&u_3\\v_1&v_2&v_3\end{vmatrix}=\left(\begin{vmatrix}u_2&u_3\\v_2&v_3\end{vmatrix},\ -\begin{vmatrix}u_1&u_3\\v_1&v_3\end{vmatrix},\ \begin{vmatrix}u_1&u_2\\v_1&v_2\end{vmatrix}\right)}

\end{fmdef}

\Needspace{7\baselineskip}
\begin{fmprop}{Propietats i aplicacions}
\begin{itemize}
\item \textbf{Anticommutativa}: \fmm{\vec v\times\vec u=-(\vec u\times\vec v)}. \fmuline{L'ordre importa!}
\item \fmm{\vec u\times\vec v=\vec 0\iff\vec u\parallel\vec v} (en particular, \fmm{\vec u\times\vec u=\vec 0}). És una manera ràpida de veure si dos vectors són para\lgem{}els o si tres punts estan alineats.
\item Distributiva: \fmm{\vec u\times(\vec v+\vec w)=\vec u\times\vec v+\vec u\times\vec w}, \fmm{\ (k\vec u)\times\vec v=k(\vec u\times\vec v)}. \textbf{No} és associativa.
\item \textbf{Vector perpendicular a dos vectors} donats: \fmm{\vec u\times\vec v}. Serà la manera habitual d'obtindre el vector normal d'un pla i el vector director d'una recta donada per dos plans.
\item \textbf{Àrees}: \fmm{\ S_{\text{paral}{\cdotp}\text{lelogram}}=|\overrightarrow{AB}\times\overrightarrow{AD}|}, \fmm{\quad S_{\text{triangle}}=\dfrac12\,|\overrightarrow{AB}\times\overrightarrow{AC}|}.
\end{itemize}
\end{fmprop}

\Needspace{7\baselineskip}
\begin{fmex}{Exemple resolt}
\begin{enumerate}[label=\textbf{\arabic*.},leftmargin=*]
\item Calcula \fmm{\vec u\times\vec v} si \fmm{\vec u=(1,2,-1)} i \fmm{\vec v=(3,0,2)}, i comprova que és perpendicular a tots dos.

\fmm{\vec u\times\vec v=\begin{vmatrix}\vec i&\vec j&\vec k\\1&2&-1\\3&0&2\end{vmatrix}=\left(\begin{vmatrix}2&-1\\0&2\end{vmatrix},\ -\begin{vmatrix}1&-1\\3&2\end{vmatrix},\ \begin{vmatrix}1&2\\3&0\end{vmatrix}\right)=(4,\ -5,\ -6)}

\fmm{(4,-5,-6)\cdot(1,2,-1)=4-10+6=0\ \checkmark\qquad(4,-5,-6)\cdot(3,0,2)=12+0-12=0\ \checkmark}
\item Calcula l'àrea del triangle de vèrtexs \fmm{A(1,0,0)}, \fmm{B(0,2,0)} i \fmm{C(0,0,3)}.

\fmm{\overrightarrow{AB}=(-1,2,0)}, \fmm{\quad\overrightarrow{AC}=(-1,0,3)}

\fmm{\overrightarrow{AB}\times\overrightarrow{AC}=\begin{vmatrix}\vec i&\vec j&\vec k\\-1&2&0\\-1&0&3\end{vmatrix}=(6,\ 3,\ 2)}, \fmm{\quad|(6,3,2)|=\sqrt{36+9+4}=7}

\fmm{S=\dfrac12\cdot7=\dfrac72\ u^2}
\end{enumerate}
\end{fmex}

\Needspace{7\baselineskip}
\section{Producte mixt}\hypertarget{sec6}{}

\Needspace{7\baselineskip}
\begin{fmdef}{Definició i propietats}
El \textbf{producte mixt} de tres vectors és el \textbf{nombre} \fmm{[\vec u,\vec v,\vec w]=\vec u\cdot(\vec v\times\vec w)}, que és igual al determinant que té els vectors per files:

\fmdisp{[\vec u,\vec v,\vec w]=\begin{vmatrix}u_1&u_2&u_3\\v_1&v_2&v_3\\w_1&w_2&w_3\end{vmatrix}}
\newpage
\begin{itemize}
\item Per les propietats dels determinants: si permutem dos vectors canvia de signe; \fmm{[\vec u,\vec v,\vec w]=[\vec v,\vec w,\vec u]=[\vec w,\vec u,\vec v]}.
\item \fmm{[\vec u,\vec v,\vec w]=0\iff} els tres vectors són \textbf{coplanaris} (linealment dependents).
\item Quatre punts \fmm{A,B,C,D} són \textbf{coplanaris} \fmm{\iff\big[\overrightarrow{AB},\overrightarrow{AC},\overrightarrow{AD}\big]=0}.
\end{itemize}
\end{fmdef}

\Needspace{7\baselineskip}
\begin{fmprop}{Volums}

\fmdisp{V_{\text{paral}{\cdotp}\text{lelepípede}}=\big|[\vec u,\vec v,\vec w]\big|\qquad\qquad V_{\text{tetraedre}\,ABCD}=\dfrac16\,\Big|\big[\overrightarrow{AB},\overrightarrow{AC},\overrightarrow{AD}\big]\Big|}

El valor absolut és important: el producte mixt pot ser negatiu, però un volum no. El tetraedre és la sisena part del para\lgem{}elepípede: el prisma triangular és la meitat del para\lgem{}elepípede i la piràmide, un terç del prisma.
\end{fmprop}

\Needspace{7\baselineskip}
\begin{fmex}{Exemple resolt}
\begin{enumerate}[label=\textbf{\arabic*.},leftmargin=*]
\item Calcula el volum del para\lgem{}elepípede determinat per \fmm{\vec u=(1,0,2)}, \fmm{\vec v=(2,1,0)} i \fmm{\vec w=(0,3,1)}.

\fmm{[\vec u,\vec v,\vec w]=\begin{vmatrix}1&0&2\\2&1&0\\0&3&1\end{vmatrix}=1\cdot(1-0)-0+2\cdot(6-0)=13}

\fmm{V=|13|=13\ u^3}
\item Calcula el volum del tetraedre de vèrtexs \fmm{A(1,0,0)}, \fmm{B(0,1,0)}, \fmm{C(0,0,1)} i \fmm{D(2,2,2)}.

\fmm{\overrightarrow{AB}=(-1,1,0)}, \fmm{\ \overrightarrow{AC}=(-1,0,1)}, \fmm{\ \overrightarrow{AD}=(1,2,2)}

\fmm{\begin{vmatrix}-1&1&0\\-1&0&1\\1&2&2\end{vmatrix}=-1\cdot(0-2)-1\cdot(-2-1)+0=2+3=5}

\fmm{V=\dfrac{|5|}{6}=\dfrac56\ u^3}
\item Troba \fmm{k} perquè els punts \fmm{A(1,0,1)}, \fmm{B(2,1,0)}, \fmm{C(0,1,2)} i \fmm{D(k,2,1)} siguen coplanaris.

\fmm{\overrightarrow{AB}=(1,1,-1)}, \fmm{\ \overrightarrow{AC}=(-1,1,1)}, \fmm{\ \overrightarrow{AD}=(k-1,2,0)}

\fmm{\begin{vmatrix}1&1&-1\\-1&1&1\\k-1&2&0\end{vmatrix}=1\cdot(0-2)-1\cdot\big(0-(k-1)\big)+(-1)\cdot\big(-2-(k-1)\big)=-2+k-1+k+1=2k-2}

\fmm{2k-2=0\Rightarrow k=1}
\end{enumerate}
\end{fmex}

\Needspace{18\baselineskip}
\section{Equacions de la recta}\hypertarget{sec7}{}

\Needspace{7\baselineskip}
\begin{fmdef}{Una recta queda determinada per un punt i un vector director}
Siga \fmm{r} la recta que passa pel punt \fmm{P(p_1,p_2,p_3)} amb \textbf{vector director} \fmm{\vec d=(d_1,d_2,d_3)\neq\vec0}. Un punt \fmm{X(x,y,z)} és de \fmm{r} si \fmm{\overrightarrow{PX}} és proporcional a \fmm{\vec d}:

\begin{center}\renewcommand{\arraystretch}{1.5}\small\begin{tabularx}{\linewidth}{>{\raggedright\arraybackslash}p{4.2cm}>{\raggedright\arraybackslash}X}\toprule
\textbf{Equació} & \textbf{Expressió}\\\midrule
Vectorial & $(x,y,z)=(p_1,p_2,p_3)+\lambda\,(d_1,d_2,d_3),\ \lambda\in\mathbb R$\\[3pt]
Paramètriques & $\begin{cases}x=p_1+\lambda d_1\\y=p_2+\lambda d_2\\z=p_3+\lambda d_3\end{cases}$\\[3pt]
Contínua & $\dfrac{x-p_1}{d_1}=\dfrac{y-p_2}{d_2}=\dfrac{z-p_3}{d_3}$\\[3pt]
Implícites (o cartesianes) & $\begin{cases}A_1x+B_1y+C_1z+D_1=0\\A_2x+B_2y+C_2z+D_2=0\end{cases}$\\[3pt]
\bottomrule\end{tabularx}\end{center}

\begin{itemize}
\item Les \textbf{equacions implícites} descriuen la recta com la intersecció de dos plans no para\lgem{}els. S'obtenen de la contínua igualant les fraccions dues a dues (multiplicant en creu).
\item \textbf{De les implícites a les paramètriques}: resolem el sistema (compatible indeterminat amb un paràmetre), o bé prenem \fmm{\vec d=\vec n_1\times\vec n_2} (producte vectorial dels vectors normals dels dos plans) i un punt qualsevol que complisca les dues equacions.
\item \textbf{Recta per dos punts} \fmm{A} i \fmm{B}: punt \fmm{A} i vector director \fmm{\overrightarrow{AB}}.
\end{itemize}
\vspace*{0.2cm}
\begin{fmnote}\fmuline{Notació}: si una component del vector director és 0, en la contínua escrivim igualment, per exemple, \fmm{\dfrac{x-1}{2}=\dfrac{y+3}{0}=\dfrac{z}{5}}. No és una divisió de veritat: indica que \fmm{y=-3} per a tots els punts de la recta.\end{fmnote}
\end{fmdef}

\Needspace{7\baselineskip}
\begin{fmex}{Exemple resolt}
\begin{enumerate}[label=\textbf{\arabic*.},leftmargin=*]
\item Escriu totes les equacions de la recta que passa per \fmm{P(1,-2,3)} amb vector director \fmm{\vec d=(2,1,-1)}.

Vectorial: \fmm{\ (x,y,z)=(1,-2,3)+\lambda(2,1,-1)}

Paramètriques: \fmm{\ \begin{cases}x=1+2\lambda\\y=-2+\lambda\\z=3-\lambda\end{cases}}\\

\vspace*{0.2cm}
Contínua: \fmm{\ \dfrac{x-1}{2}=\dfrac{y+2}{1}=\dfrac{z-3}{-1}}

Implícites: \fmm{\dfrac{x-1}{2}=\dfrac{y+2}{1}\Rightarrow x-1=2y+4\Rightarrow x-2y-5=0}; \fmm{\ \dfrac{y+2}{1}=\dfrac{z-3}{-1}\Rightarrow -y-2=z-3\Rightarrow y+z-1=0}.

\fmdisp{r:\begin{cases}x-2y-5=0\\y+z-1=0\end{cases}}

\item Escriu en forma paramètrica la recta \fmm{r:\begin{cases}x+y-z=1\\2x-y+z=2\end{cases}}

Vector director: \fmm{\ \vec n_1\times\vec n_2=\begin{vmatrix}\vec i&\vec j&\vec k\\1&1&-1\\2&-1&1\end{vmatrix}=(0,-3,-3)}, o més senzill, \fmm{\vec d=(0,1,1)}.

Un punt: fem \fmm{z=0\Rightarrow\begin{cases}x+y=1\\2x-y=2\end{cases}\Rightarrow x=1,\ y=0}. Punt \fmm{P(1,0,0)}.

\fmm{r:\begin{cases}x=1\\y=\lambda\\z=\lambda\end{cases}}. (També ho obtenim resolent el sistema amb \fmm{z=\lambda}: sumant les equacions, \fmm{3x=3}.)
\end{enumerate}
\end{fmex}

\Needspace{18\baselineskip}
\section{Equacions del pla}\hypertarget{sec8}{}

\Needspace{7\baselineskip}
\begin{fmdef}{Un pla queda determinat per un punt i dos vectors directors no para\lgem{}els}
Siga \fmm{\pi} el pla que passa per \fmm{P(p_1,p_2,p_3)} amb vectors directors \fmm{\vec u} i \fmm{\vec v} (no para\lgem{}els). Un punt \fmm{X} és de \fmm{\pi} si \fmm{\overrightarrow{PX}} és combinació lineal de \fmm{\vec u} i \fmm{\vec v}:

\begin{center}\renewcommand{\arraystretch}{1.5}\small\begin{tabularx}{\linewidth}{>{\raggedright\arraybackslash}p{3.2cm}>{\raggedright\arraybackslash}X}\toprule
\textbf{Equació} & \textbf{Expressió}\\\midrule
Vectorial & $(x,y,z)=(p_1,p_2,p_3)+\lambda\,\vec u+\mu\,\vec v,\ \ \lambda,\mu\in\mathbb R$\\[3pt]
Paramètriques & $\begin{cases}x=p_1+\lambda u_1+\mu v_1\\y=p_2+\lambda u_2+\mu v_2\\z=p_3+\lambda u_3+\mu v_3\end{cases}$\\[3pt]
General (o implícita) & $\begin{vmatrix}x-p_1&y-p_2&z-p_3\\u_1&u_2&u_3\\v_1&v_2&v_3\end{vmatrix}=0\ \Longrightarrow\ Ax+By+Cz+D=0$\\[3pt]
Segmentària & $\dfrac xa+\dfrac yb+\dfrac zc=1$ (talla els eixos en $(a,0,0)$, $(0,b,0)$, $(0,0,c)$)\\[3pt]
\bottomrule\end{tabularx}\end{center}

L'equació general diu que \fmm{\overrightarrow{PX}}, \fmm{\vec u} i \fmm{\vec v} són coplanaris (producte mixt nul).
\end{fmdef}

\Needspace{7\baselineskip}
\begin{fmprop}{El vector normal}
En l'equació general \fmm{\pi:Ax+By+Cz+D=0}, el vector \fmm{\vec n=(A,B,C)} és \textbf{perpendicular} al pla: és el \textbf{vector normal}. Per tant:

\begin{itemize}
\item \textbf{Pla per un punt amb vector normal} \fmm{\vec n}: \fmm{\ A(x-p_1)+B(y-p_2)+C(z-p_3)=0}.
\item Si coneixem dos vectors directors, \fmm{\vec n=\vec u\times\vec v}.
\item De la general a les paramètriques: aïllem una incògnita i les altres dues són els paràmetres (és un sistema d'una equació amb tres incògnites).
\end{itemize}

\textbf{Casos típics per a determinar un pla}

\begin{center}\renewcommand{\arraystretch}{1.5}\small\begin{tabularx}{\linewidth}{>{\raggedright\arraybackslash}p{7.2cm}>{\raggedright\arraybackslash}X}\toprule
\textbf{Dades} & \textbf{Punt i vectors directors}\\\midrule
Tres punts no alineats $A,B,C$ & $A;\ \overrightarrow{AB},\ \overrightarrow{AC}$\\[3pt]
Una recta $r(P,\vec d)$ i un punt $Q\notin r$ & $P;\ \vec d,\ \overrightarrow{PQ}$\\[3pt]
Dues rectes secants $r(P,\vec d_r)$ i $s(Q,\vec d_s)$ & $P;\ \vec d_r,\ \vec d_s$\\[3pt]
Dues rectes para\lgem{}eles $r(P,\vec d)$ i $s(Q,\vec d)$ & $P;\ \vec d,\ \overrightarrow{PQ}$\\[3pt]
Conté $r(P,\vec d_r)$ i és para\lgem{}el a $s$ & $P;\ \vec d_r,\ \vec d_s$\\[3pt]
Passa per $P$ i és perpendicular a la recta $r$ & $P;\ \vec n=\vec d_r$\\[3pt]
\bottomrule\end{tabularx}\end{center}
\end{fmprop}

\Needspace{7\baselineskip}
\begin{fmex}{Exemple resolt}
\begin{enumerate}[label=\textbf{\arabic*.},leftmargin=*]
\item Troba l'equació general del pla que passa per \fmm{P(1,0,2)} amb vectors directors \fmm{\vec u=(1,1,0)} i \fmm{\vec v=(0,2,-1)}.

\fmm{\begin{vmatrix}x-1&y&z-2\\1&1&0\\0&2&-1\end{vmatrix}=(x-1)(-1-0)-y(-1-0)+(z-2)(2-0)=0}

\fmm{-x+1+y+2z-4=0\Rightarrow x-y-2z+3=0}. Comprovació amb \fmm{P}: \fmm{1-0-4+3=0\ \checkmark}
\item Troba el pla que passa per \fmm{A(1,0,0)}, \fmm{B(0,2,0)} i \fmm{C(0,0,3)}.

Ja vam calcular (apartat 3.5) que \fmm{\overrightarrow{AB}\times\overrightarrow{AC}=(6,3,2)}: és un vector normal.

\fmm{6(x-1)+3y+2z=0\Rightarrow 6x+3y+2z-6=0}

Com que els punts són als eixos, també surt directament de la segmentària: \fmm{\dfrac x1+\dfrac y2+\dfrac z3=1}, i multiplicant per 6 obtenim el mateix.
\item Escriu en paramètriques el pla \fmm{x-2y+3z=4}.

Aïllem \fmm{x=4+2y-3z} i fem \fmm{y=\lambda}, \fmm{z=\mu}: \fmm{\ \begin{cases}x=4+2\lambda-3\mu\\y=\lambda\\z=\mu\end{cases}}
\end{enumerate}
\end{fmex}

\Needspace{18\baselineskip}
\section{Posicions relatives de rectes i plans}\hypertarget{sec9}{}

\Needspace{16\baselineskip}
\subsection*{A. Dues rectes}

\Needspace{7\baselineskip}
\begin{fmprop}{Amb els vectors directors i un vector que les uneix}
Siguen \fmm{r} (punt \fmm{P}, director \fmm{\vec d_r}) i \fmm{s} (punt \fmm{Q}, director \fmm{\vec d_s}). Comparem el rang de \fmm{\{\vec d_r,\vec d_s\}} amb el de \fmm{\{\vec d_r,\vec d_s,\overrightarrow{PQ}\}}:

\begin{center}\renewcommand{\arraystretch}{1.5}\small\begin{tabularx}{\linewidth}{>{\raggedright\arraybackslash}p{3.2cm}>{\raggedright\arraybackslash}X>{\raggedright\arraybackslash}X>{\raggedright\arraybackslash}X}\toprule
\textbf{$\mathrm{rg}(\vec d_r,\vec d_s)$} & \textbf{$\mathrm{rg}(\vec d_r,\vec d_s,\overrightarrow{PQ})$} & \textbf{Posició} & \textbf{Rangs amb les implícites}\\\midrule
1 & 1 & \textbf{Coincidents} & $\mathrm{rg}\,M=\mathrm{rg}\,M^*=2$\\[3pt]
1 & 2 & \textbf{Para\lgem{}eles} & $\mathrm{rg}\,M=2,\ \mathrm{rg}\,M^*=3$\\[3pt]
2 & 2 & \textbf{Secants} (es tallen en un punt) & $\mathrm{rg}\,M=\mathrm{rg}\,M^*=3$\\[3pt]
2 & 3 & \textbf{S'encreuen} (no estan en un mateix pla) & $\mathrm{rg}\,M=3,\ \mathrm{rg}\,M^*=4$\\[3pt]
\bottomrule\end{tabularx}\end{center}

A la pràctica: si \fmm{\vec d_r\parallel\vec d_s}, mirem si \fmm{P\in s} (coincidents) o no (para\lgem{}eles). Si no són para\lgem{}els, calculem \fmm{\big[\vec d_r,\vec d_s,\overrightarrow{PQ}\big]}: si és 0 són \textbf{secants} i si no, \textbf{s'encreuen}. L'última columna correspon a les rectes donades per implícites: \fmm{M} és la matriu \fmm{4\times3} dels coeficients dels quatre plans i \fmm{M^*} l'ampliada.
\end{fmprop}

\Needspace{7\baselineskip}
\begin{fmex}{Exemple resolt}
\begin{enumerate}[label=\textbf{\arabic*.},leftmargin=*]
\item Estudia la posició relativa de \fmm{r:(x,y,z)=(1,0,2)+\lambda(1,-1,2)} i \fmm{s:(x,y,z)=(2,1,0)+\mu(2,1,-1)}.

\fmm{\vec d_r=(1,-1,2)} i \fmm{\vec d_s=(2,1,-1)} no són proporcionals: les rectes es tallen o s'encreuen.

\fmm{\overrightarrow{PQ}=(1,1,-2)}, \fmm{\quad\begin{vmatrix}1&-1&2\\2&1&-1\\1&1&-2\end{vmatrix}=1\cdot(-2+1)+1\cdot(-4+1)+2\cdot(2-1)=-1-3+2=-2\neq0}

Els rangs són 2 i 3: les rectes \textbf{s'encreuen}.
\item Estudia la posició relativa de \fmm{r:(x,y,z)=(1,0,-1)+\lambda(1,2,1)} i \fmm{s:(x,y,z)=(1,3,-2)+\mu(1,-1,2)} i, si es tallen, troba el punt de tall.

\fmm{\overrightarrow{PQ}=(0,3,-1)}, \fmm{\quad\begin{vmatrix}1&2&1\\1&-1&2\\0&3&-1\end{vmatrix}=1\cdot(1-6)-2\cdot(-1-0)+1\cdot(3-0)=-5+2+3=0}

Els directors no són proporcionals i el determinant és 0: són \textbf{secants}.

\fmm{\begin{cases}1+\lambda=1+\mu\\2\lambda=3-\mu\\-1+\lambda=-2+2\mu\end{cases}\Rightarrow\lambda=\mu,\ 3\lambda=3\Rightarrow\lambda=\mu=1}. La tercera: \fmm{0=0\ \checkmark}. Punt de tall: \fmm{(2,2,0)}.
\end{enumerate}
\end{fmex}

\Needspace{16\baselineskip}
\subsection*{B. Recta i pla}

\Needspace{7\baselineskip}
\begin{fmprop}{Amb el vector normal del pla i el director de la recta}
Siguen \fmm{r} (punt \fmm{P}, director \fmm{\vec d}) i \fmm{\pi:Ax+By+Cz+D=0} (normal \fmm{\vec n}):

\begin{center}\renewcommand{\arraystretch}{2.5}\small\begin{tabularx}{\linewidth}{>{\raggedright\arraybackslash}p{3.2cm}>{\raggedright\arraybackslash}X>{\raggedright\arraybackslash}X}\toprule
\textbf{Condició} & \textbf{Posició} & \textbf{Rangs ($r$ en implícites)}\\\midrule
$\vec n\cdot\vec d\neq0$ & \textbf{Secants}: es tallen en un punt & $\mathrm{rg}\,M=\mathrm{rg}\,M^*=3$\\[3pt]
$\vec n\cdot\vec d=0,\ P\notin\pi$ & \textbf{Para\lgem{}els} & $\mathrm{rg}\,M=2,\ \mathrm{rg}\,M^*=3$\\[3pt]
$\vec n\cdot\vec d=0,\ P\in\pi$ & \textbf{La recta és continguda en el pla} & $\mathrm{rg}\,M=\mathrm{rg}\,M^*=2$\\[3pt]
\bottomrule\end{tabularx}\end{center}

Per a trobar el punt de tall, substituïm les paramètriques de la recta en l'equació del pla i aïllem \fmm{\lambda}.
\end{fmprop}

\Needspace{16\baselineskip}
\subsection*{C. Dos plans}

\Needspace{7\baselineskip}
\begin{fmprop}{Comparant els coeficients}
Siguen \fmm{\pi_1:A_1x+B_1y+C_1z+D_1=0} i \fmm{\pi_2:A_2x+B_2y+C_2z+D_2=0}:

\begin{center}\renewcommand{\arraystretch}{2.5}\small\begin{tabularx}{\linewidth}{>{\raggedright\arraybackslash}p{5.2cm}>{\raggedright\arraybackslash}X>{\raggedright\arraybackslash}X}\toprule
\textbf{Condició} & \textbf{Rangs} & \textbf{Posició}\\\midrule
$\dfrac{A_1}{A_2}=\dfrac{B_1}{B_2}=\dfrac{C_1}{C_2}=\dfrac{D_1}{D_2}$ & $\mathrm{rg}\,M=\mathrm{rg}\,M^*=1$ & \textbf{Coincidents}\\[3pt]
$\dfrac{A_1}{A_2}=\dfrac{B_1}{B_2}=\dfrac{C_1}{C_2}\neq\dfrac{D_1}{D_2}$ & $\mathrm{rg}\,M=1,\ \mathrm{rg}\,M^*=2$ & \textbf{Para\lgem{}els}\\[3pt]
$A_1,B_1,C_1$ no proporcionals a $A_2,B_2,C_2$ & $\mathrm{rg}\,M=\mathrm{rg}\,M^*=2$ & \textbf{Secants}: es tallen en una recta\\[3pt]
\bottomrule\end{tabularx}\end{center}
\end{fmprop}

\Needspace{7\baselineskip}
\begin{fmex}{Exemple resolt}
\begin{enumerate}[label=\textbf{\arabic*.},leftmargin=*]
\item Estudia la posició de \fmm{r:(x,y,z)=(1,0,2)+\lambda(2,1,-1)} i \fmm{\pi:x+y+z-6=0}.

\fmm{\vec n\cdot\vec d=(1,1,1)\cdot(2,1,-1)=2\neq0}. Són \textbf{secants}.

\fmm{(1+2\lambda)+\lambda+(2-\lambda)-6=0\Rightarrow2\lambda-3=0\Rightarrow\lambda=\dfrac32}

Punt de tall: \fmm{\ \left(1+3,\ \dfrac32,\ 2-\dfrac32\right)=\left(4,\ \dfrac32,\ \dfrac12\right)}
\item Estudia la posició dels plans \fmm{\pi_1:x-2y+z=3}, \fmm{\pi_2:2x-4y+2z=1} i \fmm{\pi_3:-2x+4y-2z+6=0}.

\fmm{\pi_1,\pi_2}: \fmm{\ \dfrac12=\dfrac{-2}{-4}=\dfrac12\neq\dfrac31\Rightarrow} \textbf{para\lgem{}els}.

\fmm{\pi_3} s'escriu \fmm{-2x+4y-2z=-6}, és a dir, \fmm{-2\cdot(x-2y+z)=-2\cdot3}: \fmm{\pi_1} i \fmm{\pi_3} són \textbf{coincidents}.

\fmuline{Compte}: per a comparar el terme independent, tots dos plans han d'estar escrits de la mateixa manera (\fmm{D} al mateix costat de la igualtat).
\end{enumerate}
\end{fmex}

\Needspace{18\baselineskip}
\section{Posició relativa de tres plans}\hypertarget{sec10}{}

\Needspace{7\baselineskip}
\begin{fmprop}{Teorema de Rouché aplicat a tres plans}
Els punts comuns de tres plans són les solucions del sistema que formen les seues equacions. Siguen \fmm{M} la matriu de coeficients i \fmm{M^*} l'ampliada:

\fmdisp{\begin{cases}A_1x+B_1y+C_1z=D_1\\A_2x+B_2y+C_2z=D_2\\A_3x+B_3y+C_3z=D_3\end{cases}\qquad M=\begin{pmatrix}A_1&B_1&C_1\\A_2&B_2&C_2\\A_3&B_3&C_3\end{pmatrix}\qquad M^*=\left(\begin{array}{ccc|c}A_1&B_1&C_1&D_1\\A_2&B_2&C_2&D_2\\A_3&B_3&C_3&D_3\end{array}\right)}

\begin{center}\renewcommand{\arraystretch}{1.5}\small\begin{tabularx}{\linewidth}{>{\raggedright\arraybackslash}p{3.2cm}>{\raggedright\arraybackslash}X>{\raggedright\arraybackslash}X>{\raggedright\arraybackslash}X}\toprule
\textbf{$\mathrm{rg}\,M$} & \textbf{$\mathrm{rg}\,M^*$} & \textbf{Sistema} & \textbf{Posició dels plans}\\\midrule
3 & 3 & SCD & \textbf{Secants en un punt}\\[3pt]
2 & 3 & SI & Cap parell de files proporcional a $M$: \textbf{secants dos a dos} (formen un prisma triangular)\\[3pt]
& & &Dues files de $M$ proporcionals: \textbf{dos plans para\lgem{}els tallats pel tercer}\\[3pt]
2 & 2 & SCI (1 paràmetre) & Cap parell de files de $M^*$ proporcional: \textbf{secants en una recta} (són d'un mateix feix)\\[3pt]
& & &Dues files de $M^*$ proporcionals: \textbf{dos plans coincidents i el tercer secant}\\[3pt]
1 & 2 & SI & Cap parell de files de $M^*$ proporcional: \textbf{tres plans para\lgem{}els}\\[3pt]
& & &Dues files de $M^*$ proporcionals: \textbf{dos plans coincidents i el tercer para\lgem{}el}\\[3pt]
1 & 1 & SCI (2 paràmetres) & \textbf{Tres plans coincidents}\\[3pt]
\bottomrule\end{tabularx}\end{center}

Al laboratori de tres plans pots veure i girar cadascun dels huit casos.
\end{fmprop}

\Needspace{7\baselineskip}
\begin{fmex}{Exemple resolt}
\begin{enumerate}[label=\textbf{\arabic*.},leftmargin=*]
\item Estudia la posició relativa de \fmm{\pi_1:x+y+z=1}, \fmm{\pi_2:x-y+2z=0} i \fmm{\pi_3:2x+3z=1}.

\fmm{|M|=\begin{vmatrix}1&1&1\\1&-1&2\\2&0&3\end{vmatrix}=1\cdot(-3)-1\cdot(3-4)+1\cdot(0+2)=-3+1+2=0}, i el menor \fmm{\begin{vmatrix}1&1\\1&-1\end{vmatrix}=-2\neq0}: \fmm{\mathrm{rg}\,M=2}.

Observem que \fmm{\pi_3=\pi_1+\pi_2} (també el terme independent: \fmm{1+0=1}). Per tant, \fmm{\mathrm{rg}\,M^*=2}.

Cap parell de plans és proporcional: els tres plans són \textbf{secants en una recta}. Si canviàrem \fmm{\pi_3} per \fmm{2x+3z=5}, tindríem \fmm{\mathrm{rg}\,M^*=3} i formarien un \textbf{prisma}.
\item Discuteix, segons els valors de \fmm{a}, la posició relativa dels plans \fmm{x+y+az=1}, \fmm{x+ay+z=1} i \fmm{ax+y+z=1}.

\fmm{|M|=\begin{vmatrix}1&1&a\\1&a&1\\a&1&1\end{vmatrix}=-(a-1)^2(a+2)\Rightarrow|M|=0\iff a=1\ o\ a=-2}

Si \fmm{a\neq1} i \fmm{a\neq-2}: \fmm{\mathrm{rg}\,M=\mathrm{rg}\,M^*=3}, \textbf{secants en un punt}.

Si \fmm{a=1}: els tres plans són \fmm{x+y+z=1}: \textbf{coincidents} (\fmm{\mathrm{rg}\,M=\mathrm{rg}\,M^*=1}).

Si \fmm{a=-2}: \fmm{\mathrm{rg}\,M=2} i el menor de \fmm{M^*} \fmm{\begin{vmatrix}1&1&1\\1&-2&1\\-2&1&1\end{vmatrix}=-9\neq0}, així que \fmm{\mathrm{rg}\,M^*=3}. Com que no hi ha dues files de \fmm{M} proporcionals, els plans són \textbf{secants dos a dos} (prisma).
\end{enumerate}
\end{fmex}

\Needspace{18\baselineskip}
\section{Feixos de plans}\hypertarget{sec11}{}

\Needspace{7\baselineskip}
\begin{fmdef}{Definicions}
\begin{itemize}
\item El \textbf{feix de plans para\lgem{}els} a \fmm{\pi:Ax+By+Cz+D=0} és el conjunt de tots els plans que tenen el mateix vector normal:

\fmdisp{Ax+By+Cz+k=0,\qquad k\in\mathbb R}

\item El \textbf{feix de plans secants} d'eix \fmm{r} és el conjunt de tots els plans que contenen la recta \fmm{r}. Si \fmm{r} ve donada per dos plans, \fmm{r:\begin{cases}\pi_1:A_1x+B_1y+C_1z+D_1=0\\\pi_2:A_2x+B_2y+C_2z+D_2=0\end{cases}}, el feix és

\fmdisp{\alpha\,(A_1x+B_1y+C_1z+D_1)+\beta\,(A_2x+B_2y+C_2z+D_2)=0,\qquad(\alpha,\beta)\neq(0,0)}

\item Normalment dividim per \fmm{\alpha} i escrivim \fmm{\pi_1+\lambda\,\pi_2=0}. \fmuline{Compte}: d'aquesta manera perdem el pla \fmm{\pi_2} (que correspon a \fmm{\alpha=0}). Si les condicions del problema no donen cap valor de \fmm{\lambda}, cal comprovar si la solució és \fmm{\pi_2}.
\end{itemize}

Els feixos són molt útils per a trobar el pla que \textbf{conté una recta donada en implícites} i compleix una altra condició: passar per un punt, ser para\lgem{}el a una recta, ser perpendicular a un pla\dots{} Imposem la condició i calculem \fmm{\lambda}.
\end{fmdef}

\Needspace{7\baselineskip}
\begin{fmex}{Exemple resolt}
\begin{enumerate}[label=\textbf{\arabic*.},leftmargin=*]
\item Troba el pla que conté la recta \fmm{r:\begin{cases}x+y-z-2=0\\2x-y+z-1=0\end{cases}} i passa per \fmm{P(1,2,3)}.

Feix de plans d'eix \fmm{r}: \fmm{\ (x+y-z-2)+\lambda\,(2x-y+z-1)=0}

Imposem que hi passe \fmm{P}: \fmm{\ (1+2-3-2)+\lambda\,(2-2+3-1)=0\Rightarrow-2+2\lambda=0\Rightarrow\lambda=1}

\fmm{(x+y-z-2)+(2x-y+z-1)=0\Rightarrow3x-3=0\Rightarrow\boxed{x=1}}
\item Troba el pla que conté la mateixa recta \fmm{r} i és para\lgem{}el a la recta \fmm{s} de vector director \fmm{\vec d_s=(1,1,0)}.

Vector normal del pla del feix: \fmm{\ \vec n=(1+2\lambda,\ 1-\lambda,\ -1+\lambda)}

Para\lgem{}el a \fmm{s\iff\vec n\perp\vec d_s}: \fmm{\ (1+2\lambda)+(1-\lambda)=0\Rightarrow\lambda=-2}

\fmm{(x+y-z-2)-2(2x-y+z-1)=0\Rightarrow-3x+3y-3z=0\Rightarrow\boxed{x-y+z=0}}
\item Troba el pla para\lgem{}el a \fmm{2x-y+3z=1} que passa per \fmm{P(1,1,1)}.

\fmm{2x-y+3z+k=0}, \fmm{\ 2-1+3+k=0\Rightarrow k=-4\Rightarrow\boxed{2x-y+3z-4=0}}
\end{enumerate}
\end{fmex}

\Needspace{16\baselineskip}
\subsection*{Recta que passa per un punt i en talla una altra}

\Needspace{7\baselineskip}
\begin{fmprop}{Com es planteja}
Per un punt \fmm{P} passen infinites rectes que tallen una recta \fmm{r} (totes estan en el pla que conté \fmm{P} i \fmm{r}). Per a fixar-ne una, l'enunciat sempre dona una \textbf{condició més}. Hi ha dos mètodes:

\begin{itemize}
\item \textbf{Punt genèric}: el punt de tall és \fmm{Q=A+\lambda\,\vec d_r\in r}. La recta buscada és \fmm{PQ}, de director \fmm{\overrightarrow{PQ}} (que depén de \fmm{\lambda}). Imposem la condició a \fmm{\overrightarrow{PQ}} i trobem \fmm{\lambda}:

\begin{itemize}
\item para\lgem{}ela a un pla \fmm{\pi}: \fmm{\ \overrightarrow{PQ}\cdot\vec n_\pi=0};
\item perpendicular a un vector \fmm{\vec u}: \fmm{\ \overrightarrow{PQ}\cdot\vec u=0} (si \fmm{\vec u=\vec d_r}, és la recta que talla \fmm{r} \textbf{perpendicularment}: ho veurem a la unitat 4);
\item que també talle una altra recta \fmm{s}: \fmm{\ \overrightarrow{PQ}}, \fmm{\vec d_s} i un vector de \fmm{P} a un punt de \fmm{s} han de ser coplanaris.
\end{itemize}
\item \textbf{Plans}: la recta buscada està en el pla \fmm{\pi_r} que conté \fmm{P} i \fmm{r} (\fmm{\vec n=\vec d_r\times\overrightarrow{AP}}, o un pla del feix d'eix \fmm{r}). Si també ha de tallar una recta \fmm{s}, està en el pla \fmm{\pi_s} que conté \fmm{P} i \fmm{s}, i és \fmm{\pi_r\cap\pi_s}. O, més ràpid: tallem \fmm{s} amb \fmm{\pi_r}; el punt de tall \fmm{Q} dona la recta \fmm{PQ}.
\end{itemize}
\vspace*{0.2cm}
\begin{fmnote}\fmuline{Compte}: al final cal comprovar que la recta \fmm{PQ} no és para\lgem{}ela a \fmm{r} (o a \fmm{s}); si ho fóra, no les tallaria i el problema no tindria solució.\end{fmnote}
\end{fmprop}

\Needspace{7\baselineskip}
\begin{fmex}{Exemple resolt}
\begin{enumerate}[label=\textbf{\arabic*.},leftmargin=*]
\item Troba la recta que passa per \fmm{P(1,0,2)}, talla la recta \fmm{r:(x,y,z)=(0,1,1)+\lambda(1,1,0)} i és para\lgem{}ela al pla \fmm{\pi:2x-y+z=0}.

Punt genèric de \fmm{r}: \fmm{\ Q(\lambda,\,1+\lambda,\,1)}, \fmm{\ \overrightarrow{PQ}=(\lambda-1,\ 1+\lambda,\ -1)}.

Para\lgem{}ela a \fmm{\pi}: \fmm{\ \overrightarrow{PQ}\cdot(2,-1,1)=2\lambda-2-1-\lambda-1=\lambda-4=0\Rightarrow\lambda=4}, \fmm{\ Q(4,5,1)}.

\fmm{\overrightarrow{PQ}=(3,5,-1)}, que no és para\lgem{}el a \fmm{\vec d_r}: \fmm{\ \boxed{(x,y,z)=(1,0,2)+t\,(3,5,-1)}}
\item Troba la recta que passa per l'origen i talla les rectes \fmm{r:(x,y,z)=(1,0,1)+\lambda(1,1,0)} i \fmm{s:(x,y,z)=(0,1,-1)+\mu(1,0,2)}.

Pla que conté \fmm{O} i \fmm{r}: \fmm{\ \vec n=\vec d_r\times\overrightarrow{OA}=(1,1,0)\times(1,0,1)=(1,-1,-1)}, \fmm{\ \pi_r:x-y-z=0}.

Tallem \fmm{s} amb \fmm{\pi_r}: \fmm{\ \mu-(1)-(-1+2\mu)=-\mu=0\Rightarrow\mu=0}, \fmm{\ Q(0,1,-1)}.

Recta \fmm{OQ}: \fmm{\ \boxed{(x,y,z)=t\,(0,1,-1)}}. Comprovació: no és para\lgem{}ela a \fmm{r}, i el punt \fmm{(0,-1,1)} (\fmm{t=-1}) és de \fmm{r} (\fmm{\lambda=-1}) .
\end{enumerate}
\end{fmex}

  
\end{document}
